\section{经济与社会：把人、地和粮接成一条战线}

一支军队从都城出发时，最先移动的往往不是将军，而是田里的收成、城门边的役夫和写在木牍上的命令。战国诸国争的当然是城池与土地，然而一块新得的土地若不能知道由谁耕种，一座新得的城若不能取得粮食、工匠和守卒，胜利便很难留在手中。于是，田界、容量、货币、道路、仓储、城垣与户口逐渐进入同一组政治问题：国家怎样把分散在乡里、市场和边地的人力物力，变成可以再次调度的资源？

这不是一条从魏、齐、楚、赵到秦整齐推进的道路。魏较早把河东、西河的谷物、徭役与军镇放在一起思考；齐的冶铁与刀币依托临淄一带的城市和东部交换；楚有江淮与中原之间的金版流通，也必须让广阔地域里的地方官文书相互衔接；赵的邯郸作坊与北方防线相连；秦则在关中田界、权量、灌溉和县邑之间加密连接。它们面对的是相似的战争压力，却各有地理、旧有关系和技术条件，不能归结为一套全国同时生效的“战国经济制度”。

\subsection{田亩、簿书与差役：国家先要看见什么}

\begin{event}{约前五世纪末至前四世纪前期}{魏、齐、赵、秦的早期物资基础}{年代与覆盖范围有争议}{河东、临淄、邯郸与关中}\eventanchor{WQ-0400-WEI-AGRICULTURE}{战国·魏的尽地力传统}\eventanchor{WQ-0398-WEI-FISCAL-GRAIN}{战国·魏的粟帛田租与徭役}\eventanchor{WQ-0400-ANYI-SPADE-MINT}{战国·安邑布币}\eventanchor{WQ-0380-QI-IRON}{战国·齐地铁器}\eventanchor{WQ-0390-IRON-CASTING-SCALE}{战国·多地铁范铸造}\eventanchor{WQ-0397-ZHAO-HANDAN-CRAFT}{战国·邯郸作坊}\eventanchor{WQ-0384-LIYANG-WALL-EVIDENCE}{战国·栎阳城垣}\eventanchor{WQ-0370-CITY-WALLS}{战国·诸国城墙}

后世将李悝同“尽地力之教”相连，魏的故事因而常被讲成战国农政的开端。这个传统至少抓住了一项真实的压力：西河驻军和中原战争不能永远依靠临时向贵族索取，谷物丰歉、田租、粟帛与徭役须成为可以反复讨论和调度的事情。可是，《汉书·食货志》和《说苑》所保存的是较晚的回望，不是一册魏国田亩簿。它们让我们知道后人把魏的强盛理解为农业、调剂与军政相接，却不能据此重建某一年、某一县的税率，或把一套后来归名于李悝的办法写成无例外的成文法。

同样，安邑及河东的空首布、平首布，显示价值可以被铸成较便于计算和转移的媒介；它并不证明田租从此全以钱收取。粟帛、实物劳役和货币能够长期并存。对魏廷来说，布币所扩大的也许是军费、市场和计量之间的选择，而不是将所有家庭突然带入同一种支付关系。河东的位置使它既接近农业腹地，也接近西河方向的军事压力；这正解释了为什么田租、服役与交换会在魏的政治记忆中彼此牵连。

齐、赵与秦看到的则是另一组物。临淄墓葬和冶铁遗址、邯郸的铜器与骨角作坊、栎阳的夯土城垣，并不替我们开出三国的产量清单，却各自留下了劳动如何在城市附近聚集的痕迹。铁范和铁制农具、兵器在多地扩展，能够增加耕作和军备的潜力，也同时需要矿料、燃料、工匠与运输；齐地的铁器较多，不能自动推出它的农具已按同一速度遍及列国。邯郸的作坊服务城邑交换和兵器需求，也不等于赵廷直接经营每一道工序。栎阳遗址跨越数代的营筑活动，说明秦在渭北经营都城和城防有可见的物质基础，却不能把每一段城墙都精确归到一位君主或一年。

城墙尤其把这些关系暴露得很清楚。诸国修筑都城和边境防御，首先是为了让粮仓、渡口、道路和居民不至于在一次突袭中失守；修城本身又要征集土石、运输和劳力。战争要求更多防御工程，工程把更多家庭和工匠带入征发，守住工程又需要更稳定的粮食与兵器供给。城垣不是“生产增长”的旁证，而是国家竞争把社会劳动固定在土地上的一种方式。
\end{event}

国家能够看见田地，并不等于它已经拥有田地，更不等于每一户的占有从此安定。秦在\eventref{WQ-0356-SHANGYANG}{战国·商鞅变法}后继续将县邑、耕战和家庭责任相连；至传统所称的\eventref{WQ-0350-QINREFORM}{战国·秦迁咸阳与第二次变法}，开阡陌、校正权衡等说法又把田界、赋役和市场计算放在一起。其意义不宜只用“土地私有化”或“废除旧制”一语概括。若田亩、粟米和劳役有较可比较的单位，县吏便较容易核验征收和差役；旧有的地权、人情和地方协商却未必因此立刻消失。

\begin{event}{约前四世纪中叶至前300年前后}{尺度、货币与地方文书进入治理}{年代与制度细节有争议}{关中、三晋、齐与楚}\eventanchor{WQ-0350-QIN-LAND-MEASURE}{战国·秦的田界与度量}\eventanchor{WQ-0344-SHANGYANG-FANGSHENG}{战国·商鞅方升}\eventanchor{WQ-0336-QIN-AGRICULTURE-MARKET}{战国·秦的农业市场管理}\eventanchor{WQ-0340-WEI-COIN}{战国·三晋布币}\eventanchor{WQ-0325-QIN-ROUND-COIN}{战国·秦地圜钱}\eventanchor{WQ-0325-QI-KNIFE-COIN}{战国·齐法化刀币}\eventanchor{WQ-0300-CHU-YINGYUAN-GOLD}{战国·楚郢爰金版}\eventanchor{WQ-0300-BAOSHAN-CHU-SLIPS}{战国·包山楚简}\eventanchor{WQ-0300-CROSSBOW-MECHANISMS}{战国·列国兵器改进}

商鞅方升是一件小器物，所能承受的解释反而格外有限。它以“十八年”标记容量标准，确实为秦的尺度实践留下了带年代的物证；一件量器却不能独自证明关中每一处市场、仓廪和田租都已完全使用同一尺度。将它同《商君书·垦令》里关于田界、垦田和市易的要求相读，更稳妥的说法是：秦的政治语言试图让土地、粮食、市场和军需可以相互核算。规范性文本说的是应当怎样治理，方升说的是有一种标准曾被铸成器物；从两者之间，仍隔着县吏、商人、农户和地方执行的距离。

三晋的布币、齐的“法化”刀币、楚的“郢爰”金版和秦地逐渐出现的圜钱、半两形制，恰好反对了“战国已有统一货币”的想象。它们的形制、材质和适用场景不同：河东的布币把魏的市场与计量联系起来，临淄的刀币同东部盐铁、谷物和手工业交换相接，楚金版适合高价值支付，秦的圆形钱制则是关中市场和行政支付的一种地方实践。钱币窖藏和器物铭文可以说明交换与计价的网络存在，不能仅由器形推出谁铸造了全部货币、政府取得多少税，或每一个农户如何买卖。

出土布币与《管子·国蓄》所保留的蓄积语言，至多让人把\eventref{WQ-0340-WEI-COIN}{战国·三晋布币}理解为市场交换和部分赋役支付的一种媒介，而不是已经覆盖所有田租的国家货币。齐的\eventref{WQ-0325-QI-KNIFE-COIN}{战国·齐法化刀币}有较整齐的铭文，秦的\eventref{WQ-0325-QIN-ROUND-COIN}{战国·秦地圜钱}则呈现另一种形制；二者提示各自的铸造与流通网络，却没有让临淄和关中共享同一套价格、税额或征收程序。对官府和居民而言，粟帛、劳役与钱币可以在同一时期承担不同的账目：前者更贴近田地与服役，后者使某些货物、军费或跨城交易较易折算。正是在这种并存而非替代中，货币把城市和乡里的关系接得更密，也把不同地区之间的制度边界保留了下来。

楚的包山简把视野从钱币压到一份份文书。诉讼、封守和书写格式让人看见楚国地方官至少在某些场合以成文程序处理事情，因而不能把楚只写成依靠封君和口头命令的“松散大国”。但一座墓中保存的简牍，只是战国中晚期某一行政片段；它不能代替楚全境的档案，更不能用来量化江汉、淮北或边远城邑的征收。楚的难题正在于地域广阔：资源丰富会增加回旋余地，却也使每一次将命令、物资和人力送往同一方向都要穿过更多地方关系。

武器的变化与这些日常计算并行。秦、楚、赵、齐所见的弩机、戈矛和较标准化的铜兵器，提示攻守战争要求零件、铜料和工坊能够持续接续；器物的产地和部队归属却并非总能判定。不能把一枚弩机的规格推成各国军队的统一装备，也不应把军事技术同农业、市场截然分开。能让兵器在需要时补上去的，是城邑作坊、道路、仓储和劳力的组合；能让农具持续出现的，同样是这些组合。
\end{event}

\begin{causalchain}[从田界到前线：一条常被遮住的因果链]
\term{结构性原因}是长期兼并使国家不能只在战时临时索取粮食、役夫和兵器。\term{使能条件}则是田界与容量可被核验、城市作坊能持续生产、市场和道路能转运不同物资。它们并不自动制造胜利，却使一次征收、一次修筑或一次调兵较可能接上下一次行动。\term{直接后果}是赋役和军需更容易进入文书与计量；\term{反馈回路}则是战争越持久，越逼迫各国强化对土地、城邑和劳力的看见。不同国家在哪一环更强、哪一环断裂，才构成了它们后来命运的差别。
\end{causalchain}

\subsection{水、路与市场：物资不会自己抵达}

\begin{event}{前286年至前242年}{市场重组、灌溉工程与前线粮道}{若干年代和规模有争议}{陶、关中、长平与东郡}\eventanchor{WQ-0286-SONG-MARKET-DISRUPTION}{战国·齐灭宋后的市场扰动}\eventanchor{WQ-0286-SONG-TAO-ROUTE}{战国·齐接收陶邑商路}\eventanchor{WQ-0260-QIN-SHANGDANG-GRAIN}{战国·长平前线转输}\eventanchor{WQ-0246-ZHENGGUO-HAN-AGENT}{战国·郑国入秦修渠}\eventanchor{WQ-0246-ZHENGGUO-CANAL-DECISION}{战国·秦续修郑国渠}\eventanchor{WQ-0246-QIN-EMPEROR-MAUSOLEUM}{战国·骊山陵园工程}\eventanchor{WQ-0242-QIN-DONGJUN-GRAIN}{战国·秦接收东郡仓田}

齐灭宋以后，陶、彭城等城邑的税源、关津和道路控制随政权易手而重组。陶地连接魏、楚、赵的商路，使齐获得西向经营的机会，也使原本依赖这些通道的商人和邻国重新估量风险。这里最不该做的是把兼并直接折算为某一方“获得多少商业利润”。传世叙事能够确认宋亡及其政治后果，货币和考古材料可以提示中原交换的密度；它们都不能替我们画出齐军到来前后货物流向的精确曲线。较可把握的是，市场与道路不是战争外面的繁荣背景，而是兼并所争夺的城邑资源的一部分。

长平把这种关系推到前线。围困赵军的秦军须由河东、上党等方向续送粟米，\eventref{WQ-0260-CHANGPING}{战国·长平之战}因而不是两位将领在孤立山谷中的对决。道路、仓储和后方征发是秦能够延长围困的使能条件，赵的守势也同样受邯郸方向粮道和补兵的限制。没有保存下来的军粮账簿，便不应给运输量、仓址或军队消耗编造整齐的数字；正因为数字不可得，反而更应把“能久战”理解为一项组织关系，而非名将个人的天赋。

郑国渠的故事把工程、阴谋和实际收益缠在一起。《史记·河渠书》将韩国水工郑国入秦说成疲秦之计；这种动机主要来自后出的叙事，不能遮蔽工程本身的水利技术价值，也不能替劳役规模提供统计。当秦王政得知郑国来历仍命续修时，较可信的历史处境是：停工会损失已投入的土石、劳力和水工知识，续修则可能扩大泾水两岸的灌溉条件。工程后来能够支持关中粮食供给，不意味着每一户都同样受益，更不意味着它没有将徭役压力推回地方。

与郑国渠相邻的骊山陵园，提醒人们不要只把大型工程视作“生产性投资”。陵园自嬴政即位后持续营建，工匠、徭役与王室礼制被长期集中于关中；其规模在不同阶段变化很大，不能把后世所见的全部遗存压回即位之年。水利、城墙、陵园和军营所争夺的，常常是同一片土地上有限的劳动力、运输能力与行政注意力。国家的工程能力固然能创造长期条件，也可能在某一时刻同家庭的耕作和生计发生正面冲突。

东郡的接收则显示秦怎样尝试把夺取变为供给。魏地仓、田与道路一旦由秦接管，东向军队便不必完全依赖临时掠取，而可以在较近的郡县节点征集和转输。可是，东郡仓储的数目、税额和每一步接收程序都没有账簿可供复原；所谓“接入网络”是一个应当逐处检验的过程，不是占城当日便全然完成的事实。秦较其他国家突出的地方，不在于它免除了后勤摩擦，而在于它较能将关中灌溉、河东转输和新占地区的仓田接成连续的战争条件。
\end{event}

\subsection{征发、迁散与城市：谁承担一场战争}

\begin{event}{前316年至前225年}{征服怎样重排城市、户口与劳力}{过程、规模与去向多有争议}{巴蜀、河东、楚北、赵魏旧地}\eventanchor{WQ-0316-SHU-MIGRATION}{战国·秦迁巴蜀军民}\eventanchor{WQ-0316-CHENGDU-WALLED-CITY}{战国·成都成为行政据点}\eventanchor{WQ-0309-QINGCHUAN-WOODEN-TABLET}{战国·青川木牍与道路整治}\eventanchor{WQ-0292-CHU-WAN-REFUGEES}{战国·楚宛地迁散}\eventanchor{WQ-0286-QIN-ANYI-RESETTLEMENT}{战国·秦迁安邑与募徙河东}\eventanchor{WQ-0248-ZHAO-TAIYUAN-REFUGEES}{战国·赵北部军民转移}\eventanchor{WQ-0234-ZHAO-BATTLEFIELD-PRISONERS}{战国·赵南部兵籍与粮食压力}\eventanchor{WQ-0228-ZHAO-HANDAN-HOUSEHOLDS}{战国·秦接收邯郸户籍}\eventanchor{WQ-0225-DALIANG-REFUGEES}{战国·大梁围灌后的迁散}

秦入巴蜀以后，迁入军民、重组城邑和设置官吏，不只是把西南涂成秦的颜色。成都的道路、城垣和可汇集粮食、文书、守军的平原位置，使它逐渐成为向南经营的节点；青川木牍所记道路整治和官府传令，则保存了山路、驿传与地方命令必须相接的一次具体痕迹。木牍的释文、地名和适用范围仍有讨论，一件文书尤其不能代表整个蜀地的行政常态。它的价值正在于提醒我们：征服后的治理并非只靠一项宏大制度，而须依靠某条道路是否整治、某道命令是否抵达的细小劳动。

河东的安邑是另一种重组。魏献安邑后，秦迁民并募徙，以补充新得地区的人力；可供确认的是迁置和招徙的政策方向，不能据此推算迁民人数，或假定所有人都按命令流动。对秦而言，城池的价值包括能够重新登记、征发和安置的人；对魏而言，失去的也不只是城墙，而是一个财政和交换的核心。城市因此既可能给居民提供市场、作坊和庇护，也可能在易主时成为户籍、差役和迁徙被重新编排的地点。

战争造成的流动并不总由胜利者主动规划。宛地失守后，楚北部居民和官吏向陈、方城等方向转移，田地、仓储与守备的旧节奏随之中断；太原方向受秦军进攻时，赵北部军民向邯郸、代地和邻近堡垒移动，地方征发被迫重新分配。两件事都没有留下可供统计的迁徙总数。不能把城市失陷说成全区人口顿时清空，也不能把避难者想成被动的数字；材料只足以让我们看见，战争将官吏、军人、家庭、粮食和市场同时推向别处，而收容它们的城邑要承担新的供给与防务。

赵在平阳、武城受挫后，战俘、伤兵和逃散者涌向邯郸周边，兵籍、医治和粮食压力一并回到都城。到\eventref{QIN-0228-ZHAO}{秦·秦破赵都}，邯郸的户籍、仓廪和工匠又由赵王税役转入秦军政管理。接收程序没有赵秦双方保存下来的完整档案，因此不能把“户籍接收”写成平滑的表格转换；但它揭示了统一战争的一项实际收益：上一国维系城市的记录和劳力，成为下一国继续北方、魏地战争的潜在资源。

大梁的结局则让水网呈现出更残酷的一面。围灌使原可用于交通、防御和生活的水系转为灾害环境，官吏、手工业者和编户迁散；具体损失和去向未有可靠统计，传世史书的渲染不能替代人口资料。秦在接收魏都时面对的，并非一座自动恢复秩序的空城，而是住房、粮食、道路和居民关系都须重新安顿的城市社会。将灭魏只写作一条军功，便会抹去水攻如何落到许多不在战场中央的人身上。
\end{event}

城市中的差异也由此显出。邯郸的冶铸、骨角作坊和边防需求，使工匠、燃料和运输成为赵国军政的一部分；临淄的冶铁、刀币和市场使齐拥有不同的东部资源；楚的金版和地方简牍所见，是高价值支付与官吏文书在广域国家中并行；关中与巴蜀的城、渠、道路和迁民，则使秦得以把不同腹地接向东进战争。这里的“社会层级”不应被写成一幅静止的等级图：传世叙事较常看见君主与将相，器物与遗址较常留下工匠、城墙和货币，简牍才偶尔保存官吏与案件。它们共同说明，不同位置的人被赋予不同的役使、生产和记录责任；却不足以替每一座城列出完整的人口、职业和收入。

\begin{event}{前221年至前216年后}{统一后的尺度、关键生产与土地核验}{年代、覆盖范围与制度解释有争议}{新建帝国各地}\eventanchor{QIN-0221-SALT-IRON-REGULATION}{秦·关键生产与尺度管理}\eventanchor{QIN-0216-LAND-TAX-BASE}{秦·自实田与地税核验}

\eventref{QIN-0230-UNIFICATION}{秦·统一六国}以后，统一尺度和官府对盐铁、兵器等关键生产交换的介入，使战国时已在不同地区形成的计量与军需问题进入更大的行政范围。不可把它写成秦在前221年立即实行毫无缝隙的全面专卖：关于官府管理的范围，尤其不能从汉代制度倒推。较稳妥的判断是，帝国需要让军械、工具、征收和交换具有较可互认的尺度，关键资源因而更容易被中央调度；地方生产与市场关系仍会保留复杂形态。

前216年后的自实田诏令同样不应被读成一把万能钥匙。它可能为地税、户籍和土地纠纷提供较清晰的核验基础，反映新帝国希望把旧国不同的田制和占有关系纳入可征发的名目；《史记》和《汉书》的相关叙事却不足以解决私有土地解释的重大分歧。国家得到更多关于田地的信息，不等于土地由此更均等，更不等于每一处乡里都稳定地接受同一口径。统一真正延续的是战国已经显出的追问：谁耕作，谁承担赋役，谁可以核验，以及一份地方记录能否被远方的权力读懂。
\end{event}

\begin{sourcenote}[同一件事，在不同材料里说的并不是同一句话]
\sourcebook{商君书}、《管子》和《墨子》提供农战、市场、蓄积、城守和器械应当如何组织的规范性语言；篇章的成书和编纂层次各有问题，不能逐句化为某国某年的执行清单。商鞅方升、布币、刀币、楚金版、铁范、城址与作坊是器物和考古材料，能说明尺度、铸造、营筑或交换曾在特定地点发生，不能由一个窖藏或一处遗址推出全国的产量、税收与生活水平。

包山楚简和青川木牍比宏大叙事更接近日常行政，却都是局部保存的文书：它们能让人看见地方官、道路和格式，不能代表楚、秦每一县的常态。《史记》《汉书》《华阳国志》及相关世家、列传保存郑国渠、迁民、围城、土地核验等历史叙事，是把事件接入长时段因果的必要线索；它们多为后出文本，尤其不能将工程规模、迁徙人数、军粮数量和税额写成已经核实的统计。把规范文本、行政简牍、考古发现和后世叙事并置，不是为了平均它们的证词，而是为了让每一种材料只回答它真正能够回答的问题。
\end{sourcenote}

\begin{debate}[“国家化”并非一张覆盖全境的平面]
铁器、灌溉、钱币、城防和户籍常被连成一条生产力与国家能力同步上升的直线。现有材料支持的是若干地区、若干环节的增强：有的地方留下冶铁和货币，有的地方留下城墙和文书，有的地方留下水利与迁民的叙事。它们的年代、扩散速度、受益群体与实际执行并不相同。因而不能从一件钱币推算市场总量，从一条渠道推算关中全部产出，也不能从一份简牍断定全国官吏同样工作。比较诸国时，重要的不是给它们排出“先进”次序，而是辨认它们把田地、城市、劳力和命令接在一起的方式为何不同。
\end{debate}

\begin{turningpoint}[制度得失：更能动员，未必更能承受]
\term{所得}：田界与权量使征收、市场和军需较可核对；城邑、道路、水利和作坊使粮食、兵器与文书较可能到达需要它们的地方；接收户籍、仓廪和工匠，则使一次征服能够成为下一次行动的条件。魏的河东经验、齐的东部市场、楚的广域文书与高值支付、赵的城市作坊和秦的关中—巴蜀—河东连接，都是这种能力的不同形态。

\term{代价与边界}：这些联系由农户的赋役、工匠的劳动、城中居民的承受和迁散者的安置维持。工程会同耕作争夺劳力，围城会破坏市场与水网，接收户籍会把旧城居民带入新的税役关系。秦后来能够把多处资源接得更紧，并不证明各地原本相同，也不保证这种高强度调度可以无限延长。国家越能将战时办法延伸到日常，越要面对家庭、地方和城市是否还能承担的限度。
\end{turningpoint}
